\title{Direct Preference Optimization: Your Language Model is Secretly a Reward Model}

\begin{abstract}
Large-scale unsupervised language models (LMs) acquire extensive world knowledge and some reasoning abilities through their training, yet exerting precise influence over their output remains challenging because their learning process lacks direct supervision. Approaches to enhance such controllability typically involve gathering human assessments of the comparative quality of model outputs and then adapting the LM to align with these preferences, often by means of reinforcement learning from human feedback (RLHF). However, RLHF is frequently intricate and unstable: it first requires constructing a reward model that captures human preferences, followed by applying reinforcement learning to optimize the unsupervised LM for this learned reward while seeking to maintain proximity to the initial model. In this work, we present a novel way to parameterize the reward model in RLHF, which allows the derivation of the associated optimal policy analytically, thereby enabling the standard RLHF task to be reformulated as a straightforward classification problem. This approach, termed \textit{Direct Preference Optimization} (DPO), is efficient, reliable, and computationally modest, removing the necessity for sampling from the LM during adaptation or extensive hyperparameter tuning. Our empirical results demonstrate that DPO aligns LMs to human preferences as well or better than existing methodologies. Importantly, DPO surpasses RLHF implementations based on PPO in controlling the sentiment of model outputs, and equals or outperforms them in generating high-quality summaries and single-turn dialogues, all while greatly reducing implementation and training complexity.
\end{abstract}

\section{Introduction}
Large-scale unsupervised language models (LMs), trained on massive datasets, have demonstrated unexpected abilities~\citep{chowdhery2022palm, brown2020language, touvron2023llama,bubeck2023sparks}. Nevertheless, these models learn from human-generated data, which reflects diverse objectives, values, and competencies. Not all of these are worth emulating; for instance, while we want an AI code assistant to \textit{recognize} prevalent coding errors so it can correct them, it is preferable for the system to favor the (sometimes infrequent) examples of expert-level coding within its data when generating code. In another example, we might expect the language model to be \textit{informed} about a misconception held by 50\% of people, but it should not reproduce that falsehood in 50\% of its responses to related prompts. Thus, determining the \emph{intended outputs and conduct} of the model—distinct from the breadth of its \textit{knowledge and skills}—is essential for developing AI that is reliable, high-performing, and controllable \citep{ouyang2022training}. Conventional strategies often use reinforcement learning (RL) to align LMs with human preferences, but as we will demonstrate, the commonly employed RL-based objective can, in fact, be directly optimized with a straightforward binary cross-entropy loss, thereby streamlining the entire preference learning process.

In broad terms, current methodologies embed desired behaviors in language models by leveraging well-crafted datasets of human preferences that represent traits considered helpful and safe. This phase of preference learning follows an initial period of extensive unsupervised pre-training on diverse text. Although supervised fine-tuning based on exemplary human responses is a direct method for preference learning, the most prominent approaches utilize reinforcement learning from human (or AI) feedback (RLHF/RLAIF; \citep{christiano2017deep,bai2022constitutional}). In RLHF, a reward model is trained on annotated human preferences and subsequently used in an RL framework to adjust the language model policy, ensuring it produces outputs associated with high reward scores while limiting deviation from the pre-trained baseline. While RLHF is effective at producing systems with advanced dialogue and coding capabilities, its learning procedure is substantially more involved than that of supervised learning, requiring the training of several models and iterative policy sampling, which greatly increases computational requirements.

This work demonstrates a method for directly aligning language models with human preferences, bypassing the need for explicit reward modeling or reinforcement learning strategies. We introduce \textit{{Direct Preference Optimization} (DPO)}, an algorithm designed to implicitly target the same goal as current RLHF approaches—maximizing rewards under a KL-divergence regularization—while remaining conceptually and practically accessible. Fundamentally, the DPO update boosts the log-odds ratio between favorable and unfavorable model responses but incorporates an adaptive, instance-specific weighting factor that guards against model collapse, a problem we observe with simplistic probability ratio-based objectives. Similar to prior methods, DPO utilizes a foundational preference model (e.g., the Bradley-Terry model; \cite{bradley1952rankanalysis}) to evaluate the alignment between reward functions and observed human preferences. Distinctly, whereas existing techniques construct a preference loss for training a reward model—then optimize a separate policy against this learned reward—DPO re-expresses the preference loss as a direct function of the policy through a change of variables. Provided with a corpus of human-labeled preferences over model outputs, DPO is able to optimize the policy itself via a straightforward binary cross-entropy loss, yielding an optimal policy corresponding to an implicit reward function consistent with the preference data.

The core contribution of this study is Direct Preference Optimization (DPO), an uncomplicated algorithm free from reinforcement learning for preference-based language model training. Our empirical findings indicate that DPO matches or surpasses the performance of prevailing methods, such as PPO-based RLHF, across tasks like sentiment transfer, summarization, and conversational modeling, for language models up to 6B parameters in size.

As self-supervised language models scale up, they increasingly demonstrate the ability to solve certain tasks in a zero-shot fashion \citep{radford2019language} or with a few examples provided in the prompt \citep{gpt3,megatron,chowdhery2022palm}. Nevertheless, direct performance on downstream applications and alignment with user expectations are substantially enhanced by subsequently fine-tuning these models on collections of instructions paired with human-generated completions \citep{mishra-etal-2022-cross,sanh2022multitask,chung2022scaling,thoppilan2022lamda}. This technique, termed `instruction-tuning,' expands the model's ability to generalize to instructions not seen during fine-tuning, thereby increasing its practical utility \citep{chung2022scaling}. Although instruction tuning is highly effective, obtaining human rankings that compare responses is often more feasible than acquiring expert demonstrations, motivating later research to adapt large language models using datasets constructed from human preferences. This strategy has led to gains in areas like translation \citep{kreutzer-etal-2018-reliability}, summarization \citep{stiennon2022learning,ziegler2020finetuning}, narrative generation \citep{ziegler2020finetuning}, and adhering to instructions \citep{ouyang2022training,ramamurthy2023is}. The standard pipeline involves training a neural reward model to predict preferences in accordance with a model such as Bradley-Terry \citep{bradley1952rankanalysis}, followed by reinforcement learning-based fine-tuning of the language model—using algorithms such as REINFORCE \citep{williams1992reinforce}, proximal policy optimization (PPO; \cite{schulman2017proximal}), or related methods \citep{ramamurthy2023is}—to maximize the reward. Parallel research harnesses instruction-tuned LLMs, additionally guided by human feedback, to create further synthetic preference datasets targeting aspects like safety and harmlessness \citep{bai2022constitutional}, often relying solely on weak human supervision through descriptive rubrics for annotation tasks. This convergence brings together approaches from reinforcement learning applied to language modeling~\citep{Ranzato2015SequenceLT,paulus2018a,wu2018learning} and from general frameworks for learning through human preferences \citep{christiano2017deep,kupcsik2018learning}. Nonetheless, applying reinforcement learning to optimize large language models for human preferences poses significant practical difficulties; this work proposes a theoretically-grounded alternative to achieve preference optimization without relying on RL.

Beyond linguistic applications, the challenge of learning policies based on preferences has received significant attention in both the bandit and reinforcement learning literature, yielding a variety of proposed strategies. In the case of contextual bandit algorithms, learning is driven by action preferences or rankings as opposed to explicit rewards; this framework is commonly referred to as the contextual dueling bandit (CDB; \cite{yue2012karmed,dudik2015contextual}). Without access to definitive reward signals, theoretical analyses of CDBs employ the concept of a \textit{von Neumann winner}, defined as a policy whose expected probability of outperforming any alternative policy meets or exceeds 50\% \citep{dudik2015contextual}. Notably, preference feedback in CDBs is provided in an online fashion, whereas in the human preference learning domain, policy learning generally proceeds from a predetermined, static batch of action pairs annotated with preference data \citep{yan2022human}. Analogously, in \textit{preference-based RL} (PbRL), agents are trained on binary preference comparisons produced by an unobserved 'scoring' function, as opposed to direct reward signals \citep{BusaFekete2014,ruiz2023dueling}. Numerous PbRL techniques have been introduced, with some specifically tailored to leverage off-policy preference datasets. However, these methods typically require first estimating the underlying scoring function (i.e., constructing a reward model), which is then optimized \citep{jain2013learning,BusaFekete2014,christiano2017deep,sadigh2017active,kupcsik2018learning}. By contrast, our work proposes a unified policy learning paradigm, wherein policy parameters are optimized directly to align with demonstrated preferences, bypassing separate reward model estimation.

\section{Preliminaries}\label{section:prelims}

We summarize the RLHF procedure described by \citeauthor{ziegler2020finetuning} (as well as later adaptations such as \citep{stiennon2022learning, bai2022training, ouyang2022training}). This methodology commonly comprises three core steps: (1) supervised fine-tuning (SFT); (2) preference collection and reward function construction; and (3) reinforcement learning-driven optimization.

\textbf{SFT}: The RLHF workflow typically initiates with supervised fine-tuning of a pre-trained language model, leveraging high-quality annotated datasets tailored to the target application (such as dialogue or summarization), yielding a refined model $\pi^\text{SFT}$.

\textbf{Reward Modelling Phase}: Subsequently, the fine-tuned SFT model is supplied with prompts $x$ to generate answer pairs $(y_1, y_2)\sim \pi^\text{SFT}(y \mid x)$. Human annotators are then presented with these answer pairs and are tasked with indicating a preference for one completion, recorded as $y_w\succ y_l \mid x$—where $y_w$ denotes the answer chosen as preferable and $y_l$ the less favored among $(y_1, y_2)$. It is postulated that these preferences arise from an underlying, unobservable reward function $r^*(y, x)$. Several statistical models can be employed to interpret preference data, with the Bradley-Terry (BT) model \cite{bradley1952rankanalysis} frequently utilized. Broader ranking models such as the Plackett-Luce framework \citep{plackett1975analysis, luce2012individual} can also be applied in scenarios where more than two responses are ranked. In the BT formulation, the human-generated preference probabilities $p^*$ are modeled as follows:

\begin{equation}\label{eq:bradley-terry}
    p^*(y_1\succ y_2 \mid x)=\frac{\exp\left(r^*(x, y_1)\right)}{\exp\left(r^*(x, y_1)\right) + \exp\left(r^*(x, y_2)\right)}.
\end{equation}

Given a fixed set of comparison samples $\mathcal{D}=\bigl\{x^{(i)}, y_w^{(i)}, y_l^{(i)}\bigr\}_{i=1}^N$ drawn from $p^*$, a reward model $r_{\phi}(x, y)$ can be defined and its parameters determined via maximum likelihood estimation. By casting the preference prediction task as binary classification, the loss corresponds to the negative log-likelihood:

\begin{equation}\label{eq:reward_model}
    \mathcal{L}_R(r_{\phi}, \mathcal{D}) = -\mathbb{E}_{(x, y_w, y_l)\sim \mathcal{D}}\bigl[\log \sigma(r_{\phi}(x, y_w)- r_{\phi}(x, y_l))\bigr]
\end{equation}

where $\sigma$ denotes the logistic sigmoid. In large language model applications, $r_{\phi}(x, y)$ typically reuses the architecture and initialization of the SFT model $\pi^\text{SFT}(y \mid x)$, but with a final linear head producing a scalar score, as in \cite{ziegler2020finetuning}. To stabilize reward learning, prior studies often normalize the reward values so that $\mathbb{E}_{x,y\sim \mathcal{D}}\left[r_\phi(x, y)\right] = 0$ for all $x$.

\textbf{RL Fine-Tuning Phase}: At the RL fine-tuning stage, the trained reward model provides supervision to the policy. In accordance with previous research~\citep{jaques2017sequence, jaques2020human}, optimization is formulated as:

\begin{equation}\label{eq:RL}
\max_{\pi_{\theta}}  \mathbb{E}_{x\sim \mathcal{D}, y\sim \pi_{\theta}(y \mid x)}\bigl[r_{\phi}(x, y)\bigr] - \beta\mathbb{D}_{\textrm{KL}}\bigl[\pi_{\theta}(y\mid x)\mid \mid \pi_\text{ref}(y\mid x)\bigr],
\end{equation}

with $\beta$ tuning the strength of regularization with respect to the reference policy $\pi_\text{ref}$, namely $\pi^\text{SFT}$. Typically, the policy $\pi_\theta$ also inherits its initialization from $\pi^\text{SFT}$. The KL term is critical for ensuring the policy remains close to regions where the reward model is trustworthy, and for promoting response diversity while discouraging collapse to deterministic high-reward outputs. Because generating language is inherently discrete and non-differentiable, this problem is generally addressed with reinforcement learning. The prevailing method \citep{ziegler2020finetuning, stiennon2022learning, bai2022training, ouyang2022training} expresses the reward as ${r(x, y) = r_{\phi}(x, y) -\beta (\log \pi_{\theta}(y\mid x) - \log \pi_\text{ref}(y\mid x))}$, and optimizes it with PPO \cite{schulman2017proximal}.

\section{Direct Preference Optimization}\label{sec:DPO}

Recognizing the computational burdens that arise when applying reinforcement learning methods to large-scale model fine-tuning, we aim to propose an alternative, streamlined approach that operates directly on preference data. Instead of the conventional two-stage RLHF paradigm—reward modeling followed by RL-based policy learning—our method constructs a specific reward function parameterization allowing for closed-form retrieval of the associated optimal policy, thus circumventing the need for an explicit RL phase. The central contribution is an analytic transformation relating reward models and their optimal policies, facilitating direct conversion of reward-based loss objectives into policy-based objectives. Through this reparameterization, explicit reward network fitting is unnecessary, yet the learning process continues to reflect prevailing models of human choice such as the Bradley-Terry framework. Consequently, the policy architecture simultaneously embodies both the language generation mechanism and the implicitly encoded reward function.

\textbf{Deducing the DPO objective.} Our derivation begins with the RL objective defined in Eq.~\ref{eq:RL}, which is parametrized by a generic reward function $r$. As established in previous studies~\citep{peters2007reinforcement, peng2019advantage, korbak2022reinforcement, go2023aligning}, one can directly verify that the optimal policy for the KL-regularized reward maximization problem in Eq.~\ref{eq:RL} is given by:

\begin{equation}\label{eq:op_policy}
    \pi_r(y\mid x) = \frac{1}{Z(x)}\pi_\text{ref}(y\mid x)\exp\left(\frac{1}{\beta}r(x, y)\right),
\end{equation}

where $Z(x) =\sum_{y}\pi_\text{ref}(y\mid x)\exp\left(\frac{1}{\beta}r(x, y)\right)$ serves as the normalizing constant or partition function. The full derivation is provided in Appendix \ref{app:derivation1}. Direct computation of $Z(x)$ can be computationally challenging, even if the reward function is approximated using the MLE estimator $r_{\phi}$ in place of the true reward $r^*$~\citep{korbak2022reinforcement, go2023aligning}. This hinders the practicality of the form in Eq.~\ref{eq:op_policy}. Nevertheless, Eq.~\ref{eq:op_policy} can be rearranged to articulate the reward as a function of the optimal policy $\pi_r$, the baseline policy $\pi_\text{ref}$, and $Z(\cdot)$. By taking the logarithm of both sides of Eq.~\ref{eq:op_policy} and performing algebraic manipulation, we arrive at:

\begin{equation}\label{eq:main_eq}
    r(x,y) =\beta \log \frac{\pi_r(y\mid x)}{\pi_\text{ref}(y\mid x)} + \beta \log Z(x).
\end{equation}

This transformation can be specifically applied to the true reward $r^*$ and its optimal policy $\pi^*$. The Bradley-Terry preference model considers only reward differences between pairs of outputs, namely, ${p^*(y_1 \succ y_2 \mid x) = \sigma(r^*(x, y_1) - r^*(x, y_2))}$. Upon substituting Eq.~\ref{eq:main_eq} for $r^*(x,y)$ into the preference formulation (Eq.~\ref{eq:bradley-terry}), the partition function $Z(x)$ cancels out, thus allowing the probability of human preference to depend solely on $\pi^*$ and $\pi_\text{ref}$. Therefore, under the Bradley-Terry model, the optimal RLHF policy $\pi^*$ fulfills the preference relation:

\begin{equation}\label{eq:objective}
    p^*(y_1\succ y_2 \mid x)=\frac{1}{1 + \exp\left(\beta \log \frac{\pi^*(y_2\mid x)}{\pi_\text{ref}(y_2\mid x)} - \beta \log \frac{\pi^*(y_1\mid x)}{\pi_\text{ref}(y_1\mid x)}\right)}
\end{equation}

Appendix~\ref{app:derivation2} contains a comprehensive derivation. While Eq.~\ref{eq:objective} is built on the Bradley-Terry model, analogous results follow for more general choice models such as the Plackett-Luce model~\citep{plackett1975analysis, luce2012individual}; these are provided in Appendix~\ref{app:plackett_luce_models}.

Having now expressed the likelihood of human preference data in terms of the optimal policy (instead of the reward function), we can recast the problem as maximum likelihood estimation for a parametrized policy $\pi_\theta$. In close analogy to the reward-based objective (see Eq.~\ref{eq:reward_model}), the corresponding policy learning objective can be written as:

\begin{equation}\label{eq:optimum_model}
    \mathcal{L}_\text{DPO}(\pi_{\theta}; \pi_\text{ref}) = -\mathbb{E}_{(x, y_w, y_l)\sim \mathcal{D}}\left[\log \sigma \left(\beta \log \frac{\pi_{\theta}(y_w\mid x)}{\

In this formulation, we estimate an implicit reward through an alternative form of parameterization, in which the resulting optimal policy is simply $\pi_\theta$. Additionally, because our method aligns with fitting a reparameterized Bradley-Terry model, it inherits established theoretical guarantees—such as consistency, assuming appropriate conditions on the distribution of the preference data \cite{bong2022generalized}. Section~\ref{sec:theory} further elaborates on the theoretical aspects of DPO and compares them with prior research.

\textbf{How does the DPO update work?} To gain a mechanistic perspective on DPO, we examine the gradient of its loss, $\mathcal{L}_\text{DPO}$. The derivative with respect to $\theta$ is given by:

\begin{multline*}\label{eq:gradient}
    \nabla_\theta \mathcal{L}_\text{DPO}(\pi_\theta;\pi_\text{ref}) = \\ -\beta\mathbb{E}_{(x, y_w, y_l) \sim \mathcal{D}} \bigg[\underbrace{\sigma(\hat{r}_\theta(x, y_l) - \hat{r}_\theta (x, y_w))}_\text{assigns greater emphasis when the reward estimation errs}\bigg[\underbrace{\nabla_\theta\log \pi(y_w \mid x)}_\text{favor $y_w$ by raising its probability} - \underbrace{\nabla_\theta\log\pi(y_l \mid x)}_\text{penalize $y_l$ by reducing its probability}\bigg]\bigg],
\end{multline*}

where $\hat{r}_\theta(x, y) = \beta \log \frac{\pi_\theta(y \mid x)}{\pi_\text{ref}(y \mid x)}$ represents the reward inferred from the language model $\pi_\theta$ in comparison to the reference model $\pi_\text{ref}$ (see Section~\ref{sec:theory} for further details). Conceptually, the gradient of $\mathcal{L}_\text{DPO}$ acts to make the preferred outputs $y_w$ more likely while simultaneously making the disfavored completions $y_l$ less probable. A key aspect is that the contribution of each data point is scaled according to the degree to which the implicit reward model $\hat{r}_\theta$ incorrectly prefers $y_l$ over $y_w$, adjusted by $\beta$. This reflects the extent of error in ranking according to the reward model and incorporates the impact of the KL regularization strength. Our empirical results underscore the necessity of this weighting; omitting the scaling factor leads to suboptimal outcomes, sometimes causing language model collapse (see Appendix Table~\ref{tab:unlikelihood_generations}).

\textbf{DPO overview.} The typical DPO workflow includes: 1) Drawing completions $y_1, y_2 \sim \pi_\text{ref}(\cdot \mid x)$ for each prompt $x$, assigning human preference labels, and constructing an offline preference dataset $\mathcal{D} = \{x^{(i)}, y_w^{(i)}, y_l^{(i)}\}_{i=1}^N$; and 2) training the language model $\pi_\theta$ by minimizing $\mathcal{L}_\text{DPO}$ given $\pi_\text{ref}$, $\mathcal{D}$, and target $\beta$. 
Practically, it is preferable to leverage existing public preference datasets rather than collect new samples and annotations. Since these datasets are generally created using $\pi^\text{SFT}$, we set $\pi_\text{ref} = \pi^\text{SFT}$ when possible. If $\pi^\text{SFT}$ is unavailable, we set $\pi_\text{ref}$ to maximize the likelihood over preferred responses ${(x, y_w)}$, i.e., ${\pi_\text{ref} = \argmax_{\pi}\mathbb{E}_{x, y_w \sim \mathcal{D}}\left[\log \pi(y_w \mid x)\right]}$. This strategy reduces the discrepancy between the inaccessible true reference distribution and the practical $\pi_\text{ref}$ in DPO training. Further details regarding implementation specifics and selected hyperparameters are provided in Appendix~\ref{app:implementation}.

\section{Theoretical Analysis of DPO} In this section, we further analyze DPO, offer theoretical justification, and connect its strengths to limitations observed in RLHF actor-critic methods (such as PPO~\cite{schulman2017proximal}).

\label{sec:theory}

\subsection{Your Language Model Is Secretly a Reward Model} DPO sidesteps both explicit reward model fitting and reinforcement learning for policy optimization by employing a unified maximum likelihood formulation. The optimization problem in Eq. \ref{eq:main_eq} aligns with a Bradley-Terry model parameterized by the reward $r^*(x, y) = \beta \log\frac{\pi^*_\theta(y \mid x)}{\pi_\text{ref}(y \mid x)}$. Training the parameterized model $\pi_{\theta}$ corresponds—under an appropriate change of variables—to reward model learning as formulated in Eq. \ref{eq:reward_model}. This section establishes the theoretical foundation for this reparameterization, demonstrates that it imposes no limitations on the class of reward models learnable, and that it enables exact retrieval of the optimal policy. We begin by introducing an equivalence relation among reward functions.

\begin{definition}
Two reward functions $r(x, y)$ and $r'(x, y)$ are said to be equivalent if and only if there exists a function $f$ such that $r(x, y) - r'(x, y) = f(x)$.     
\end{definition}

This equivalence relation is readily verified, and as a result, it partitions the space of reward functions into distinct equivalence classes. We next present two key lemmas:

\begin{lemma}\label{lemma:same_prefrence}
Within the Plackett-Luce framework, and specifically in the Bradley-Terry model, any two reward functions that belong to the same equivalence class will generate identical preference distributions.
\end{lemma}

\begin{lemma}\label{lemma:same_policy}
    If two reward functions are from the same equivalence class, they induce identical optimal policies in the associated constrained RL problem.
\end{lemma}

The arguments for these results are elementary and are therefore postponed to Appendix \ref{app:lemma1}. The first lemma highlights a classical under-specification property inherent in models from the Plackett-Luce family \cite{plackett1975analysis}. Because of this ambiguity, it is standard practice to introduce identifiability constraints to secure meaningful maximum likelihood estimates according to Eq. \ref{eq:reward_model} \cite{bong2022generalized}. The second lemma establishes that any reward function from a given equivalence class leads to the same optimal policy, which implies that, for the purposes of optimization, it suffices to recover any representative reward function from the appropriate class. In Appendix~\ref{app:thm1}, we establish the following result:

\begin{theorem}\label{thm:main}
    Provided certain regularity conditions hold, every reward equivalence class consistent with the Plackett-Luce (and, in particular, Bradley-Terry) models admits a representative of the form ${r(x, y) = \beta \log \frac{\pi(y\mid x)}{\pi_\text{ref}(y\mid x)}}$ for some model $\pi(y\mid x)$ and some specified reference model $\pi_\text{ref}(y \mid x)$.
\end{theorem}

\begin{sproof}
    Let $r(x, y)$ denote any reward function, and let $\pi_r(y \mid x)$ be its corresponding optimal model, as defined in Eq. \ref{eq:op_policy}. We demonstrate that within the equivalence class of $r$, there exists a reward function expressible by the aforementioned reparameterization. Define the following projection:
\begin{equation}
    f(r; \pi_\text{ref}, \beta)(x, y) = r(x, y) - \beta\log\sum_{y}\pi_\text{ref}(y\mid x)\exp\left(\frac{1}{\beta}r(x, y)\right)
\end{equation}
Here, the operator $f$ rescales the reward by subtracting the (logarithm of the) normalization factor, which depends only on the prefix $x$. Consequently, $f(r; \pi_\text{ref}, \beta)(x, y)$ remains within the equivalence class of $r(x, y)$. Substituting the reward expression from Eq.~\ref{eq:main_eq} (which universally applies), one obtains $f(r; \pi_\text{ref}, \beta)(x, y) = \beta \log \frac{\pi_r(y\mid x)}{\pi_\text{ref}(y\mid x)}$. Therefore, the operator $f$ yields an equivalence-class member matching the desired representation, confirming that no generality is lost when using the proposed reparameterization for reward functions.
\end{sproof}

An alternative interpretation of Theorem~\ref{thm:main} is that it precisely identifies the reward function, within each equivalence class, that the DPO reparameterization chooses. Specifically, it selects the reward function that satisfies:

\begin{equation}\label{eq:lag_p}
     \sum_{y}\underbrace{\pi_\text{ref}(y\mid x)\exp\left(\frac{1}{\beta}r(x, y)\right)}_{=\pi(y\mid x)\text{, using Thm.~\ref{thm:main} reparam.}} = 1,
\end{equation}

which means that $\pi(y\mid x)$ defines a legitimate probability distribution: it assigns positive probabilities that sum to one. Notably, Eq.~\ref{eq:lag_p}—in light of Eq.~\ref{eq:op_policy}—is the partition function corresponding to the optimal policy generated by the reward function $r(x, y)$. The central insight underpinning the DPO algorithm is its ability to enforce particular restrictions on the otherwise under-determined Plackett-Luce family of preference models (including Bradley-Terry models), ensuring the retained representational power for reward models while making the optimal policy from Eq. \ref{eq:op_policy} computationally manageable across all prompt contexts $x$.

\subsection{Instability of Actor-Critic Algorithms}

Our analytical framework also illuminates sources of instability encountered with standard actor-critic algorithms commonly applied to RLHF tasks, such as PPO. Focusing on the RL fine-tuning stage as outlined in Section \ref{section:prelims}, we establish a connection to the control-as-inference perspective~\cite{levine2018reinforcement} in the context of the constrained RL formulation given by \ref{eq:RL}. Here, considering a parametric policy $\pi_{\theta}(y\mid x)$, we aim to minimize the divergence $\mathbb{D}_{\text{KL}}[\pi_{\theta}(y|x) \mid \mid \pi^*(y\mid x)]$, with $\pi^*$ as the optimal policy from Eq. \ref{eq:optimum_model}, which is itself defined by the reward function $r_{\phi}(y, x)$. Through manipulation, the resulting objective can be expressed as follows:

\begin{equation}\label{eq:AC}
    \max_{\pi_{\theta}}\mathbb{E}_{\pi_{\theta}(y\mid x)}\bigg[\underbrace{r_{\phi}(x, y) -\beta\log\sum_{y}\pi_\text{ref}(y\mid x)\exp\left(\frac{1}{\beta}r_{\phi}(x, y)\right)}_{f(r_{\phi}, \pi_\text{ref}, \beta)} - \underbrace{\beta\log\frac{\pi_{\theta}(y\mid x)}{\pi_\text{ref}(y\mid x)}}_{\text{KL}}\bigg]
\end{equation}

The objective under consideration aligns with that targeted in earlier studies \citep{ziegler2020finetuning, stiennon2022learning, bai2022training, ouyang2022training}, utilizing the DPO-equivalent reward within the reward function family $r_{\phi}$. Here, the normalization component in $f(r_{\phi}, \pi_\text{ref}, \beta)$ can be understood as representing the soft value function of the reference policy $\pi_\text{ref}$. Though this normalization does not influence the set of optimal solutions, omitting it can cause the policy gradient estimator for the objective to have excessive variance, leading to instability during optimization. One approach is to substitute this normalization term with a trainable value function; however, this introduces further optimization challenges. Alternatively, previous research has managed reward normalization by leveraging a human completion as a baseline, essentially implementing a single-sample Monte Carlo approximation of the normalization. The DPO reparameterization, by contrast, produces a reward function inherently free from the need for any baseline adjustments.

\section{Experiments}
This section presents empirical investigations into the effectiveness of DPO for training policies based solely on preference data. Our initial analysis, situated within a controlled text generation context, examines how DPO manages the trade-off between reward maximization and minimizing KL-divergence from the reference policy, in comparison to established preference optimization methods such as PPO. We then extend our evaluation to encompass larger model architectures and more challenging RLHF scenarios, including summarization and conversational tasks. Remarkably, DPO achieves performance on par with or superior to robust baselines like RLHF with PPO or the selection of the highest-rewarded trajectory among $N$ samples, all while requiring minimal hyperparameter tuning. Prior to detailing these empirical outcomes, we summarize our experimental methodology; comprehensive experimental specifications can be found in Appendix~\ref{app:exp_details}.

\textbf{Tasks.} We conduct experiments on three distinct open-ended text generation tasks. In each task, the methods learn a policy using a dataset of preferences $\mathcal{D}=\bigl\{x^{(i)}, y_w^{(i)}, y_l^{(i)}\bigr\}_{i=1}^N$. For \textbf{controlled sentiment generation}, the input $x$ consists of the beginning of a movie review sourced from the IMDb dataset \cite{maas-EtAl:2011:ACL-HLT2011}, and the objective for the policy is to generate a continuation $y$ exhibiting positive sentiment. To ensure controlled assessment, this experiment utilizes preference pairs automatically generated through a pre-trained sentiment classifier, selecting pairs where $p(\text{positive}\mid x,y_w)>p(\text{positive}\mid x,y_l)$. The SFT approach involves fine-tuning GPT-2-large until the model converges, using review data from the IMDb training set (see App~\ref{app:sentiment_details} for additional information). In the \textbf{summarization} task, $x$ is a Reddit forum post, with the policy required to generate a summary $y$ that distills the key points. Consistent with previous studies, the Reddit TL;DR summarization dataset \citep{volske-etal-2017-tl} is employed, together with human-labeled preferences from \citeauthor{stiennon2022learning}. The SFT model is obtained by further fine-tuning on a set of human-crafted Reddit summaries\footnote{\url{https://huggingface.co/CarperAI/openai_summarize_tldr_sft}}, using the TRLX framework \citep{leandro_von_werra_2023_7790115} for RLHF. Human preference data, as collected by \citeauthor{stiennon2022learning}, derive from a different yet comparably fine-tuned SFT model. For the \textbf{single-turn dialogue} setting, $x$ is a user-issued prompt, which may cover topics ranging from astrophysics to interpersonal advice. Here, the policy's task is to craft a helpful and engaging reply $y$. We make use of the Anthropic Helpful and Harmless dialogue dataset \citep{bai2022training}, which comprises 170,000 conversational exchanges between humans and an automated assistant. Each entry concludes with two assistant-generated responses, plus a label that indicates human preference. Since there is no existing SFT model for this case, we fine-tune a publicly available language model solely on the preferred responses to establish the SFT model.

\textbf{Evaluation.} We adopt two primary strategies for evaluation in our experiments. To assess the performance of each algorithm in maximizing constrained rewards, we utilize the controlled sentiment generation context, where the reward frontier—mapping achieved reward against KL-divergence from the reference policy—can be explicitly computed due to our access to the true reward signal (provided by a sentiment classifier). In contrast, outside this controlled scenario, where the genuine reward function is unavailable, we rely on a \textit{win rate} metric measured relative to a baseline policy. Specifically, we use GPT-4 to approximate human judgment on summary quality and response helpfulness in summarization and single-turn dialogue, respectively. For the summarization task, the baseline consists of reference summaries from the test set; for dialogue, it is the response indicated as preferred in the dataset. Prior work indicates language models may serve as more effective automatic evaluators than traditional metrics \citep{Chen2023ExploringTU}, but we additionally corroborate our use of GPT-4 through a human subject study, detailed in Sec.~\ref{sec:human-judgments}. Results demonstrate that GPT-4’s assessments align closely with those of humans, with agreement rates comparable to, or surpassing, inter-annotator agreement among humans.

\textbf{Methods.} Alongside DPO, we examine various prevalent methodologies for aligning language model behavior with human preferences. The simplest approaches include zero-shot prompting of \textbf{GPT-J} \citep{gpt-j} for summarization and two-shot prompting of \textbf{Pythia-2.8B} \citep{biderman2023pythia} for dialogue. We further consider the \textbf{SFT} model as well as \textbf{Preferred-FT}, which entails supervised finetuning on the preferred outputs $y_w$—these are derived from SFT (for sentiment and summarization) or from a standard language model (for single-turn dialogue). Another related approach is the \textbf{Unlikelihood} method~\citep{welleck2019neural}, which modifies training objectives so as to boost the probability of $y_w$ and explicitly decrease the likelihood of $y_l$, the latter effect modulated by an adjustable hyperparameter $\alpha\in[0,1]$. Additionally, we investigate \textbf{PPO} \citep{schulman2017proximal} models using reward signals inferred from preference data, and \textbf{PPO-GT}, which acts as an oracle by learning directly from the true reward function available in the controlled sentiment setting. For sentiment experiments, two variants of PPO-GT are employed: an off-the-shelf version \cite{leandro_von_werra_2023_7790115}, and a custom implementation featuring reward normalization and tailored hyperparameter tuning for enhanced outcomes (these modifications are likewise used for PPO models with learned rewards). As a further point of comparison, the \textbf{Best of $N$} approach involves sampling $N$ candidate responses from the SFT model (or Preferred-FT for dialogue), then selecting the response with the highest score according to the preference-trained reward model. While this strategy achieves strong results and isolates reward model quality from PPO training dynamics, it becomes computationally burdensome at moderate $N$, as it requires generating $N$ outputs per test query.

\subsection{How well can DPO optimize the RLHF objective?}

The standard KL-constrained reward maximization framework in RLHF algorithms is designed to maximize reward without allowing the policy to diverge excessively from a designated reference policy. Thus, fair algorithm comparison requires consideration of both the attained reward and the KL divergence; securing marginally greater rewards at the cost of substantially higher KL divergence is not always advantageous. In Figure~\ref{fig:frontier-tldr-main}, the relationship between reward and KL for several algorithms is depicted within the sentiment classification scenario. For each method, we conduct several independent training sessions, varying the hyperparameter that modulates policy conservativeness in each case (using target KL $\in\{3,6,9,12\}$ for PPO, $\beta \in \{0.05,0.1,1,5\}$, $\alpha\in\{0.05,0.1,0.5,1\}$ for unlikelihood, and diverse random seeds for preferred-FT). In total, 22 distinct training instances are included in this exploration. At intervals of 100 training steps until the model converges, every policy is evaluated against a collection of test prompts; during this evaluation, we calculate both the mean reward according to the true reward metric and the mean sequence-level KL\footnote{Specifically, this is computed as the aggregate of per-timestep KL-divergences.} with respect to the reference policy, $\text{KL}\left(\pi\mid \mid \pi_\text{ref}\right)$. Our analysis indicates that DPO establishes the most favorable frontier, attaining superior rewards while also maintaining a lower KL compared to alternatives. This outcome is striking for several reasons. Foremost, both DPO and PPO are optimized under the same objective, yet DPO demonstrates greater efficiency: DPO consistently yields a reward/KL tradeoff that is superior to PPO. Furthermore, DPO's frontier outperforms that of PPO, \emph{even in situations where PPO is provided with access to true rewards} (PPO-GT).

\subsection{Can DPO scale to real preference datasets?}
\label{sec:dpo-real-datasets}
We proceed to assess the effectiveness of DPO when fine-tuned on real-world preference datasets for tasks such as summarization and single-turn dialogue. In the case of summarization, automatic metrics like ROUGE have been shown to align poorly with human judgments~\citep{stiennon2022learning}, and previous studies have demonstrated that PPO-based fine-tuning on human preference data can yield more desirable summaries. For our analysis, we generate model outputs on the test portion of the TL;DR summarization dataset and calculate the average rate at which these outputs outperform reference summaries from the test set. We evaluate samples from each method across a range of temperatures (0.0 to 1.0), presenting the win rate results in Figure~\ref{fig:frontier-tldr-main} (right). DPO, PPO, and Preferred-FT all start from the same GPT-J SFT initialization\footnote{\url{https://huggingface.co/CarperAI/openai_summarize_tldr_sft}}. Our findings indicate that DPO attains an approximate win rate of 61\% at a temperature setting of 0.0, surpassing PPO’s top win rate of about 57\% at its optimal (also 0.0) sampling temperature. Furthermore, DPO achieves a superior peak win rate when compared to the $N$-best baseline method. It is important to point out that the $\beta$ parameter in DPO was not extensively optimized, suggesting that there may be further room for performance gains. Additionally, we observe that DPO demonstrates greater resilience to changes in sampling temperature compared to PPO, whose performance tends to regress to the level of the original GPT-J at higher temperatures. The Preferred-FT approach shows minimal gains relative to the base SFT model. In Section~\ref{sec:human-judgments}, we further provide human evaluation results directly comparing DPO and PPO, where outputs from DPO at a temperature of 0.25 were favored in 58\% of cases relative to PPO outputs at temperature 0.

For single-turn dialogue experiments, we assess various approaches using the portion of the Anthropic HH dataset's test split \citep{bai2022training} that contains a single round of interaction between a human and an assistant. To calculate win rates across methods, GPT-4's evaluation uses the test set’s preferred completions as the reference standard. In the absence of a standard SFT baseline for this setting, we initialize with a pre-trained Pythia-2.8B model and employ Preferred-FT on the selected completions to produce a reference model, ensuring generated outputs remain within the model’s distribution. Subsequent training is carried out using DPO. Our comparisons also include the strongest result from 128 Preferred-FT completions (as the Best of $N$ method yields diminishing returns beyond $N=128$; refer to Appendix Figure~\ref{fig:best-of-n}) and a two-shot prompting strategy with the Pythia-2.8B base model; we find that DPO matches or surpasses the best outcomes for each method at their optimal temperature settings. Additionally, we test an RLHF model trained with PPO on the Anthropic HH dataset\footnote{\url{https://huggingface.co/reciprocate/ppo_hh_pythia-6B}} and sourced from a reputable implementation\footnote{\url{https://github.com/CarperAI/trlx/tree/main/examples/hh}}, but are unable to identify a combination of prompts or sampling temperatures that outperforms the original Pythia-2.8B model. Drawing on our findings from TL;DR and the shared reward objective of the two techniques, we regard the Best of 128 strategy as an approximate stand-in for PPO-level effectiveness. Taken together, DPO emerges as the sole method that is both computationally practical and capable of improving over the preferred completions in the Anthropic HH dataset, and its results are competitive with or superior to those of the much more resource-intensive Best of 128 baseline. Moreover, as demonstrated in Figure~\ref{fig:dialogue-main}, DPO reaches peak performance within relatively few training steps.

\subsection{Generalization to a new input distribution}

To assess how PPO and DPO perform when confronted with a distribution shift, we test the policies learned from our Reddit TL;DR summarization study on a distinct dataset: news articles from the test partition of the CNN/DailyMail corpus \citep{nallapati-etal-2016-abstractive}. For this evaluation, we utilize the optimal sampling temperatures identified on TL;DR (0 and 0.25). The outcomes are shown in Table~\ref{tab:ood}. We measure the GPT-4 win rate against the reference summaries from these datasets, employing the same GPT-4 (C) prompt as in the Reddit TL;DR task, but substituting ``forum post'' for ``news article''. On this new input domain, the DPO policy consistently surpasses the PPO policy by a notable margin. These findings suggest that DPO-trained policies are capable of generalizing comparably well to those obtained by PPO, despite the fact that DPO is not exposed to the additional unlabeled Reddit TL;DR prompts available to PPO.

\subsection{Validating GPT-4 judgments with human judgments}
\label{sec:human-judgments}
We conduct a human evaluation to check the consistency and validity of GPT-4’s assessments, leveraging outputs from the TL;DR summarization task and employing two variants of GPT-4 prompts. The \textbf{GPT-4 (S)} (simple) prompt requests a selection of the summary that better captures the crucial content of the post. The \textbf{GPT-4 (C)} (concise) prompt in addition asks which summary is more succinct; this is tested as we observe that GPT-4, under the \textbf{GPT-4 (S)} prompt, shows a preference for lengthier, sometimes redundant summaries, in contrast to human raters. The complete set of prompts can be found in Appendix~\ref{app:prompts}. Our evaluation comprises three pairwise comparisons: the highest-performing system (DPO, temp. 0.25), the lowest-performing system (PPO, temp. 1.0), and a mid-level system (SFT, temp. 0.25), each contrasted with the PPO model at its best sampling temperature (greedy). This approach is intended to represent a spectrum of summary qualities; for DPO and PPO-1, we collect multiple human ratings due to limited availability of annotators. Analysis reveals that, for both prompts, the agreement between GPT-4 and humans is comparable to inter-annotator human agreement, indicating that GPT-4 can serve as a reliable proxy for human judgment. Among the two prompts, \textbf{GPT-4 (C)} generally yields win rates that more closely reflect human preferences, and is therefore chosen for reporting our main results in Section~\ref{sec:dpo-real-datasets}. More comprehensive details on the human evaluation procedure, including the interface shown to raters and a roster of annotators, are provided in Appendix~\ref{app:human-study}.

\section{Discussion}
Learning from preferences represents an effective and extensible approach for training language models that are both competent and aligned with human values. In this work, we have presented DPO, a straightforward methodology for training language models using preference data without relying on reinforcement learning. Rather than transforming the preference learning objective into a conventional reinforcement learning formulation to leverage established RL algorithms, DPO establishes a correspondence between policies of language models and reward functions. This correspondence makes it possible to train language models to align with human preferences \textit{directly} using a simple cross-entropy loss function, thus eliminating the need for reinforcement learning or any loss of generality. With minimal adjustment of hyperparameters, DPO achieves performance comparable to, and often surpassing, prevalent RLHF algorithms, such as those utilizing PPO; consequently, DPO significantly lowers the barrier to training language models based on human preferences.

\textbf{Limitations \& Future Work.} Several open questions emerge from our findings that warrant further investigation. How well does the policy learned by DPO extrapolate to out-of-distribution scenarios, in contrast to approaches trained on explicit reward signals? Preliminary evidence indicates that DPO-trained models generalize comparably to those developed with PPO, but systematic and rigorous analyses are still needed. For instance, is it feasible for DPO to effectively exploit unlabeled prompts through self-labeling techniques? Furthermore, it remains to be explored how reward over-optimization may occur in the direct preference optimization framework, and whether the slight dip in performance illustrated in Figure~\ref{fig:dialogue-main}-right can be attributed to this phenomenon. Moreover, our present evaluation is confined to models with up to 6B parameters; advancing DPO to operate efficiently at the scale of current state-of-the-art, much larger models is a promising research avenue. In terms of assessment protocols, we observe that win rates measured by GPT-4 are sensitive to prompt variations, suggesting that further work is required to identify the optimal procedures for obtaining reliable judgments from automated evaluators. Lastly, the utility of DPO extends well beyond its use in aligning language models, with significant potential for application to training generative models in additional modalities.